\begin{frame}{에피소드가 끝나면: \texttt{terminated}와
\texttt{truncated}를 구분}
  둘 다 ``이 에피소드는 여기까지''지만 \textbf{끝난 이유가 다름}.
  \begin{columns}[T]
    \begin{column}{0.5\textwidth}
      \kb{\texttt{terminated} = True}
      \begin{itemize}
        \item \textbf{환경 자체}가 ``게임이 끝났다''고 선언.
        \item CartPole: 막대 12도 초과, 카트 레일 이탈.
        \item 터미널 상태 --- 그 뒤의 미래 가치
          $V(s_{t+1}) = 0$ (게임이 끝났으니까).
      \end{itemize}
    \end{column}
    \begin{column}{0.5\textwidth}
      \kb{\texttt{truncated} = True}
      \begin{itemize}
        \item 환경이 끝나지 않았으나 \textbf{사람이 정한 한도}(최대
          스텝 수) 도달로 \textbf{강제} 종료. CartPole: 500 스텝.
        \item 게임은 \textbf{계속됐을 것} --- 다음 상태의 가치는
          그대로 $V(s_{t+1})$을 써야.
      \end{itemize}
    \end{column}
  \end{columns}
  \vskip0.6em
  \begin{block}{왜 나누는가}
    \textbf{가치함수를 재귀적으로 업데이트}할 때(5$\sim$7주차 MC/TD)
    이 둘이 다름. 에피소드 종료 판단은
    $\texttt{done} = \texttt{terminated} \text{ or }
    \texttt{truncated}$로 하되, \textbf{가치 업데이트}에서는
    \texttt{terminated}만 ``미래 가치 0'' 신호로 쓰는 것이 정석.
    구분하지 않으면 \textbf{500스텝째에 Q값을 0으로 잘라} 학습이
    부정확해짐.
  \end{block}
\end{frame}
